\section{Agentic Thinking 的本质}

\subsection{不同的优化目标}

Agentic thinking 与 reasoning thinking 有着根本不同的优化目标：

\begin{itemize}
    \item \textbf{Reasoning thinking}：由最终答案前的内部推理质量来评判——能否解出定理、写出证明、产出正确代码、通过 benchmark
    \item \textbf{Agentic thinking}：由在与环境交互过程中能否持续取得进展来评判
\end{itemize}

\begin{importantbox}{核心问题的转变}
中心问题从 ``Can the model think long enough?''（模型能否思考足够久？）转变为 ``Can the model think in a way that sustains effective action?''（模型能否以维持有效行动的方式来思考？）
\end{importantbox}

\subsection{Agentic Thinking 的独特挑战}

Agentic thinking 必须处理 reasoning 模型可以基本回避的几个问题：

\begin{enumerate}
    \item \textbf{行动时机判断}：决定何时停止思考并采取行动
    \item \textbf{工具选择与排序}：选择调用哪个工具以及以什么顺序
    \item \textbf{噪声观测整合}：整合来自环境的噪声或部分观测
    \item \textbf{失败后修正}：在失败后修订计划
    \item \textbf{长程一致性}：在多轮对话和多次工具调用中维持连贯性
\end{enumerate}

作者用一句话精炼概括：\textit{Agentic thinking is a model that reasons through action.}（Agentic thinking 是通过行动来推理的模型。）

\subsection{本章小结}

Agentic thinking 的核心不是``思考更久''，而是``以维持有效行动的方式思考''。它引入了行动时机、工具编排、噪声整合、失败恢复和长程一致性等 reasoning 模型无需面对的挑战。


